% This template was initially provided by Dulip Withanage.
% Modifications for the database systems research group
% were made by Conny Junghans,  Jannik Strötgen and Michael Gertz

\documentclass[
     12pt,         % font size
     a4paper,      % paper format
     %BCOR=10mm,     % binding correction
     DIV=14,        % stripe size for margin calculation
%     liststotoc,   % table listing in toc
%     bibtotoc,     % bibliography in toc
%     idxtotoc,     % index in toc
%     parskip       % paragraph skip instad of paragraph indent
     ]{scrreprt}

%%%%%%%%%%%%%%%%%%%%%%%%%%%%%%%%%%%%%%%%%%%%%%%%%%%%%%%%%%%%

\usepackage[parfill]{parskip}
\usepackage{amsmath,capt-of,lipsum, amssymb, bbold}
\usepackage{amsthm}
\usepackage{algorithm}
\usepackage{algpseudocode}
\usepackage{natbib}
\usepackage{graphicx}
\usepackage{caption}
\usepackage{tikz}
% \usepackage{subfig}
\usepackage{subcaption}
\usepackage[section]{placeins}
% PACKAGES:

% Use German :
\usepackage[english]{babel}
% Input and font encoding
\usepackage[utf8]{inputenc}
\usepackage[T1]{fontenc}
% Index-generation
\usepackage{makeidx}
% Einbinden von URLs:
\usepackage{url}
% Special \LaTex symbols (e.g. \BibTeX):
%\usepackage{doc}
% Include Graphic-files:
\usepackage{graphicx}
% Include doc++ generated tex-files:
%\usepackage{docxx}
% Include PDF links
%\usepackage[pdftex, bookmarks=true]{hyperref}
\usepackage{setspace}
\usepackage{array}
% hyperrefs in the documents
\usepackage[bookmarks=true,colorlinks,pdfpagelabels,pdfstartview = FitH,bookmarksopen = true,bookmarksnumbered = true,linkcolor = black,plainpages = false,hypertexnames = false,citecolor = black,urlcolor=black]{hyperref} 
%\usepackage{hyperref}

\usepackage{varwidth}
\DeclareCaptionFormat{myformat}{%
  % #1: label (e.g. "Table 1")
  % #2: separator (e.g. ": ")
  % #3: caption text
  \begin{varwidth}{\linewidth}%
    \centering
    #1#2#3%
  \end{varwidth}%
}

\usepackage{url}

%%%%%%%%%%%%%%%%%%%%%%%%%%%%%%%%%%%%%%%%%%%%%%%%%%%%%%%%%%%%

% OTHER SETTINGS:

% Pagestyle:
\pagestyle{headings}

% Choose language
\newcommand{\setlang}[1]{\selectlanguage{#1}\nonfrenchspacing}
\newcolumntype{P}[1]{>{\centering\arraybackslash}p{#1}}

\newtheorem{theorem}{Theorem}[section]
\newtheorem{definition}{Definition}[section]


\usepackage{mathtools}
\DeclarePairedDelimiter\ceil{\lceil}{\rceil}
\DeclarePairedDelimiter\floor{\lfloor}{\rfloor}

\usepackage{hyperref}

\date{\today}

\begin{document}

% TITLE:
\pagenumbering{roman} 
\begin{titlepage}


\vspace*{1cm}
\begin{center}
\vspace*{3cm}
\textbf{ 
\Large University of Heidelberg\\
\smallskip
\Large Institute of Computer Science\\
\smallskip
}

\vspace{2cm}

\textbf{\Large Project Report\\[0.4 cm]	 \large Advanced Machine Learning} % Bachelor-Arbeit 

\vspace{1.5\baselineskip}
\rule{\linewidth}{0.2 mm} \\[0.4 cm]
{\huge
\textbf{Weakly-Supervised Surgical Tool}\\
\bigskip
\textbf{Detection and Localization}
}\\
\bigskip
\rule{\linewidth}{0.2 mm} \\[1.0 cm]

by\\[0.5cm]


\begin{table}[h]
	\large
	\centering
	\begin{tabular}{P{5.5cm}P{5cm}}
		Fabian Schneider & Konrad Goldenbaum \\ 
		4017026 & 42069 \\
	\end{tabular}
\end{table}

\vspace{0.5cm}
\includegraphics[width=0.4cm]{github}\hspace{0.15cm}\href{https://github.com/faberno/SurgicalToolLocalization}{faberno/SurgicalToolLocalization}
\vfill

\today


\end{center}

\end{titlepage}

\chapter*{Abstract}
Surgical Tool Localization is a discipline with lots of useful applications. In traditional surgeries it could be used to evaluate the work of a surgeon and to find possible mistakes, while it could also be applied to the vast amount of already existing surgical videos, that then in return could be used to train even more complex systems like Surgical Workflow Analyzers. Intuitively, such a system would be trained fully supervised with lots of spatially annotated image or video data, which however is infeasible because of the sheer amount of images that would first have to be processed by humans and least be checked by additional professionals. Instead we will show that it only takes binary tool presence labels - which are relatively easy to generate, even for lots of data - to train a Fully Convolutional Neural Network to accurately localize surgical tools. We solved the problem using a pretrained convolutional neural network as a backbone and a convolutional layer as the head of the network, in order to localize objects on a heat-map.

\chapter*{Notation}
\input{0_abstract/notation}

\newpage

% MAIN PART:
% Table of contents (dmke)
\tableofcontents
\cleardoublepage
\pagenumbering{arabic} 

% List of figures (Abbildungsverzeichnis):
%\listoffigures
% List of tables (Tabellenverzeichnis):
%\listoftables

%%%%%%%%%%%%%%%%%%%%%%%%%%%%%%%%%%%%%%%%%%%%%%%%%%%%%%%%%%%%%%%
% Here, the actual content of your thesis begins
% You can either put all the text here or use individual files to store the chapters of your thesis.
% Below are templates for both alternatives.

\chapter{Introduction}
\label{chap:intro}
what is the problem? why is it interesting? what are the main obstacles? (2 Seiten)

\section{Motivation}
Reliable surgical tool localization is a milestone on the way to applications that enable assessments of surgical performance and efficiency, identifications of the skillful use of instruments and choreographies as well as planning of operative and logistical aspects of surgical resources. In addition, the ability to automatically detect and track surgical instruments in endoscopic videos paves the way for much more complex problems such as computer-assisted or even fully automated surgery. This could foster advances in the field of telemedicine, for example, by taking on simpler parts of surgeries fully automatically, or by making the automatic delivery of instruments more fluid. These advances are only possible with the help of large data sets indicating not only the presence but also the location of instruments. This is a difficult task because it is tedious and time consuming to obtain the annotations needed to train machine learning models. This is because the bounding boxes in videos have to be labelled frame by frame so far, and for a large number of surgical tools and operations. In addition, annotators need to be trained to keep up with innovations in surgical instruments. Automated localization and classification of surgical tools makes it possible to create higher-quality datasets cost-effectively by allowing a model to learn to classify and localize instruments based on a weakly annotated dataset. Weakly annotated in this context means that for each clip only the presence of the respective tools is given.

\section{Objectives \& Structure}
\subsection{Objectives}

Our goal was to implement a ML-model-architecture that can be trained on different weakly-annotated endoscopic image-datasets to label and localize tools. 

\subsection{Structure}
In a first part we will describe the problem in more depth. After that we will give a short introduction into existing literature, both on image-classification as well as localization tasks, with a focus on papers that cover similar problems like the one at hand. Connected to this introduction into existing literature we cover ML-tools that should be known in order to better follow the actual contend of the work. 
In the methodology-chapter we firstly dive into the early considerations and present our first approach, since this allows us to better illustrate the challenges and discoveries made in the process. Additionally we describe the data-management-process.
We then explain the model we finally developed and the training-strategy. We finish naming our quality metrics and success-definition.
The perhaps most illustrating part, the experiments-section, helps explain the design-decisions we made on the way. Additionally, we describe in greater depth the three datasets we used. We close encouraging further work in the field.

\newpage


\chapter{Background and Related Work}
\label{chap:background}

detailed description of the problem, discussion of prior work, introduction of relevant machine learning tools (2 Seiten)

\chapter{Methodology}
\label{chap:methodology}

detailed description of your own approach, discussion of alternatives, introduction of the experimental methodology and quality metrics (7-8 Seiten)

\chapter{Experiments}
\label{chap:experiments}

description of experiments, diagrams, comparison with results of prior work, interpretation of the experimental outcomes (7-8 Seiten)

\chapter{Summary and Future Work}
\label{chap:summary}
1 Seite


\chapter{Appendix}
\label{chap:appendix}
Beschreibung vom Code (2 Seiten)










\bibliographystyle{apalike}
% b) The File:
\bibliography{references}

\end{document}
